\documentclass[11pt]{article}

\usepackage[letterpaper,margin=1in]{geometry}
\usepackage{amsmath,amssymb}
\usepackage{booktabs}
\usepackage{longtable}
\usepackage{array}
\usepackage{xcolor}
\usepackage{listings}
\usepackage{url}
\usepackage{hyperref}
\usepackage{titlesec}
\usepackage{parskip}
\usepackage{fancyhdr}
\usepackage{enumitem}

% Hyperref colors
\hypersetup{
  colorlinks=true,
  linkcolor=black,
  urlcolor=blue!50!black,
  citecolor=black,
  pdftitle={KV Cache Engine --- Interface Specification},
  pdfauthor={LonghornSilicon},
  pdfsubject={ISA Specification, Version kv-isa-0.2}
}

% Headings
\titleformat{\section}{\normalfont\Large\bfseries}{\thesection.}{0.6em}{}
\titleformat{\subsection}{\normalfont\large\bfseries}{\thesubsection}{0.6em}{}
\titleformat{\subsubsection}{\normalfont\normalsize\bfseries}{\thesubsubsection}{0.6em}{}

% Headers / footers
\pagestyle{fancy}
\fancyhf{}
\lhead{\small KV Cache Engine ISA}
\rhead{\small \texttt{kv-isa-0.2}}
\cfoot{\thepage}
\renewcommand{\headrulewidth}{0.4pt}

% Code listings
\definecolor{codebg}{gray}{0.97}
\definecolor{codekey}{rgb}{0.13,0.29,0.53}
\definecolor{codecom}{rgb}{0.25,0.5,0.25}
\lstset{
  basicstyle=\ttfamily\small,
  backgroundcolor=\color{codebg},
  frame=single,
  framerule=0.4pt,
  rulecolor=\color{black!30},
  xleftmargin=0.5em,
  xrightmargin=0.5em,
  aboveskip=0.6em,
  belowskip=0.6em,
  showstringspaces=false,
  breaklines=true,
  keywordstyle=\color{codekey}\bfseries,
  commentstyle=\color{codecom}\itshape,
  numbers=none
}

\setlist{itemsep=1pt,topsep=2pt}

\begin{document}

%==============================================================================
% Title page
%==============================================================================
\thispagestyle{empty}
\begin{center}
\vspace*{1.2in}
{\Huge\bfseries KV Cache Engine}\\[0.4em]
{\LARGE\bfseries Interface Specification}\\[1.2em]

{\large \textsc{LonghornSilicon LLM Inference Accelerator}}\\[0.4em]
{\large Block 2 of 4 --- KV Cache Compression / Decompression}\\[2em]

\rule{0.7\textwidth}{0.4pt}\\[1em]

\begin{tabular}{rl}
\textbf{Version}:    & \texttt{kv-isa-0.2} (ChannelQuant codec) \\
\textbf{Date}:       & 2026-07-03 \\
\textbf{Status}:     & Codec integrated; RTL signed off (all CI gates green, Sky130) \\
\textbf{Author}:     & Chaithu Talasila, The University of Texas at Austin \\
\textbf{Reference}:  & \texttt{LonghornSilicon/kv-cache-engine} \\
\end{tabular}\\[1em]
\rule{0.7\textwidth}{0.4pt}

\vfill
\begin{minipage}{0.85\textwidth}
\small
\textbf{Scope.} This document describes the externally-visible interface to the
\texttt{kv\_cache\_engine} block as it will appear on the chip, on an
FPGA prototype, or in the bit-accurate software model.
A compiler backend targeting the KV cache engine writes against this
interface; the hardware implementation must conform to it. \textbf{v0.2} replaces
the retired TurboQuant+ codec (PolarQuant + QJL + Walsh--Hadamard rotation) with
\textbf{ChannelQuant} --- per-channel INT4 keys, per-token INT4 values, and static
outlier-channel isolation --- following the algorithm contract frozen in the
sibling \texttt{channelquant} repository. This specification is stable for the
KV cache engine block only and will be unified with the rest of the LonghornSilicon
ISA when the Token Importance Unit and Memory Hierarchy Controller blocks land.
\end{minipage}

\vspace{0.3in}
\end{center}

\clearpage

%==============================================================================
% TOC
%==============================================================================
\tableofcontents
\clearpage

%==============================================================================
% 1. Block overview
%==============================================================================
\section{Block Overview}
\label{sec:overview}

The KV cache engine manages on-chip storage of key and value tensors for
transformer attention, with lossy compression to reduce DRAM bandwidth. It
implements \textbf{ChannelQuant}, which quantizes on the axis that dominates the
KV error budget for grouped-query-attention (GQA) models --- a few fixed
high-magnitude \emph{key channels}:

\begin{itemize}
\item \textbf{Key path (per-channel INT4)}: keys are scaled \emph{per channel}
      over a group of \texttt{G} tokens, then quantized to INT4. The top-$k$
      worst channels (tier CQ-4+) are held FP16 instead of INT4.
\item \textbf{Value path (per-token INT4)}: values are scaled per token and
      quantized to INT4 (INT8 for the CQ-8 tier).
\end{itemize}

\noindent\textbf{Tiers} (\texttt{INFO\_TIER}):
\begin{itemize}
\item \textbf{CQ-8}: per-token INT8 for both K and V.
\item \textbf{CQ-4}: per-channel INT4 keys, per-token INT4 values.
\item \textbf{CQ-4+}: CQ-4 plus $k$ FP16 outlier key channels from a static
      calibrated ROM mask.
\end{itemize}

\noindent The compression flow is:
\begin{enumerate}
\item \textbf{Key group buffer}: accumulate \texttt{G} key tokens
      (\texttt{residual\_buffer}).
\item \textbf{Per-channel scale}: take the per-channel max over the group
      (\texttt{amax\_unit}, key mode) and freeze \texttt{D} per-channel FP16
      scales (\texttt{scale\_bank}).
\item \textbf{Quantize}: keep channels $\to$ INT4; outlier channels (CQ-4+) $\to$
      FP16 (static ROM mask). Values quantize per token against a per-token FP16
      scale.
\item \textbf{Store}: a unified per-channel SRAM record (\S\ref{sec:algorithm}).
\end{enumerate}

\noindent Decompression is per channel: \texttt{code $\cdot$ scale} reconstructs
each keep channel; outlier channels replay their stored FP16 directly. The
read-side bus is \textbf{IEEE-754 fp32} (contract \S1).

\medskip
\noindent\textbf{Compression:} $\approx$4 effective bits/value, $\approx$3.8$\times$
vs.\ FP16, near-lossless.\\
\textbf{Accuracy (HellaSwag acc\_norm, n=1000):} Qwen2-0.5B 0.421/0.415 (CQ-4/CQ-4+)
vs.\ 0.426 FP16; Qwen2-1.5B 0.505/0.517 vs.\ 0.522 FP16 (CQ-4+ helps at D=128, not D=64).\\
\textbf{Bit-exact reference:} \path{sw/reference_model/channelquant_ref.{hpp,cpp}}
and the frozen Python reference in \path{../channelquant/reference/}.

%==============================================================================
% 2. Address space and memory map
%==============================================================================
\section{Address Space and Memory Map}
\label{sec:memmap}

The block exposes a 256-byte AXI-Lite slave register window. All registers
are 32-bit, word-aligned. The base address is set at chip integration time.

\begin{table}[h]
\centering
\small
\caption{AXI-Lite register window (offsets relative to base address), ISA v0.2.}
\label{tab:regs}
\begin{tabular}{@{}l l c l p{6cm}@{}}
\toprule
Offset & Name & Access & Reset & Purpose \\
\midrule
\texttt{0x00} & \texttt{CTRL}              & RW    & 0x0   & bit[0]: \texttt{soft\_reset} (write-1 to pulse); bit[1]: \texttt{enable} \\
\texttt{0x04} & \texttt{STATUS}            & R     & 0x1   & bit[0]: \texttt{idle}; bit[3]: \texttt{sram\_full} \\
\texttt{0x08} & \texttt{INFO\_DIM}         & R     & (syn) & Head dimension \texttt{VECTOR\_DIM} ($D$) \\
\texttt{0x0C} & \texttt{INFO\_TIER}        & R     & (syn) & 0 = CQ-8, 1 = CQ-4, 2 = CQ-4+ \\
\texttt{0x10} & \texttt{INFO\_GROUP}       & R     & (syn) & Key group size \texttt{G} (contract \S3.1) \\
\texttt{0x14} & \texttt{INFO\_SRAM\_DEPTH} & R     & (syn) & Number of SRAM entries \\
\texttt{0x18} & \texttt{INFO\_CR\_K}       & R     & (syn) & Key compression ratio, fixed-point 8.8 \\
\texttt{0x1C} & \texttt{INFO\_CR\_V}       & R     & (syn) & Value compression ratio, fixed-point 8.8 \\
\texttt{0x20} & \texttt{INFO\_VERSION}     & R     & 0x00020000 & bits[31:24] major, [23:16] minor, [15:8] patch, [7:0] build \\
\texttt{0x24} & \texttt{OCCUPANCY}         & R     & 0x0   & Number of valid entries in SRAM \\
\texttt{0x28} & \texttt{WRITE\_ADDR}       & RW    & 0x0   & Target write / key-group-base address \\
\texttt{0x2C} & \texttt{READ\_ADDR}        & RW    & 0x0   & Target read address (write launches a decompress) \\
\texttt{0x30} & \texttt{KV\_SELECT}        & RW    & 0x0   & bit[0]: 0 = key, 1 = value \\
\texttt{0x34} & \texttt{IRQ\_MASK}         & RW    & 0x0   & bits[3:0]: interrupt enable mask \\
\texttt{0x38} & \texttt{IRQ\_STATUS}       & R/W1C & 0x0   & bits[3:0]: interrupt pending; write-1-to-clear \\
\texttt{0x3C} & \texttt{INFO\_OUTLIER\_K}  & R     & (syn) & Top-$k$ FP16 outlier key channels (CQ-4+) \\
\texttt{0x40} & \texttt{INFO\_SCALE\_DEPTH}& R     & (syn) & Per-channel scale-bank depth ($=D$) \\
\texttt{0x44} & \texttt{INFO\_RESID\_DEPTH}& R     & (syn) & Residual-buffer depth ($=G$) \\
\bottomrule
\end{tabular}
\end{table}

\medskip
\noindent\textbf{Conventions.} \texttt{RW} = read/write; \texttt{R} = read-only;
\texttt{W1C} = write-1-to-clear; (syn) = value fixed at synthesis time.
Writes to read-only registers are silently dropped. Reading reserved offsets
returns \texttt{0xDEADBEEF}.

%==============================================================================
% 3. Streaming data interfaces
%==============================================================================
\section{Streaming Data Interfaces}
\label{sec:stream}

Two AXI-Stream interfaces carry the vector data. The AXI-Lite window in
\S\ref{sec:memmap} is for control only.

\subsection{\texttt{s\_axis\_kv} --- Vector Write (Slave)}
\label{sec:s_axis_kv}

\begin{table}[h]
\centering
\small
\begin{tabular}{@{}l c c p{8cm}@{}}
\toprule
Signal      & Width                  & Direction & Purpose \\
\midrule
\texttt{tdata}  & \texttt{COORD\_WIDTH} (16) & in    & One FP16 coordinate \\
\texttt{tlast}  & 1                     & in        & Assert on the final coordinate of the token \\
\texttt{tvalid} & 1                     & in        & Handshake: data is valid \\
\texttt{tready} & 1                     & out       & Handshake: block is ready to accept \\
\texttt{tuser}  & 1                     & in        & 0 = key vector, 1 = value vector \\
\bottomrule
\end{tabular}
\end{table}

\noindent\textbf{Protocol:}
\begin{itemize}
\item A complete token is exactly \texttt{VECTOR\_DIM} beats. \texttt{tlast}
  asserts on the final beat.
\item \texttt{tuser} selects the path: value tokens (and, in the CQ-8 tier, key
  tokens) compress per token; key tokens in CQ-4/CQ-4+ enter the \emph{grouped}
  per-channel path.
\item \textbf{Key grouping.} Program \texttt{WRITE\_ADDR} to the group base, then
  stream \texttt{G} key tokens back-to-back. The engine buffers the group, freezes
  the per-channel scales, and emits the group's per-token records to
  \texttt{base}, \texttt{base+1}, \dots. (Partial-group flush, $g<G$, is a
  planned extension; the datapath already supports it.)
\item Backpressure: \texttt{tready} deasserts during compression / SRAM write /
  group emit. Upstream must hold signals stable.
\end{itemize}

\subsection{\texttt{m\_axis\_kv} --- Vector Read (Master)}
\label{sec:m_axis_kv}

\begin{table}[h]
\centering
\small
\begin{tabular}{@{}l c c p{8cm}@{}}
\toprule
Signal      & Width                    & Direction & Purpose \\
\midrule
\texttt{tdata}  & \texttt{OUT\_WIDTH} (32) & out    & One IEEE-754 fp32 decompressed coordinate \\
\texttt{tlast}  & 1                       & out       & Asserted on the final coordinate \\
\texttt{tvalid} & 1                       & out       & Handshake: data is valid \\
\texttt{tready} & 1                       & in        & Handshake: downstream is ready to accept \\
\bottomrule
\end{tabular}
\end{table}

\noindent\textbf{Protocol.} Writing \texttt{READ\_ADDR} launches a decompress; the
block then streams \texttt{VECTOR\_DIM} fp32 beats, one channel per beat. The
read bus is fp32 (contract \S1); a downstream attention core may narrow to fp16
outside the codec's parity boundary.

\subsection{Eviction Interface}
\label{sec:evict}

\begin{table}[h]
\centering
\small
\begin{tabular}{@{}l c c p{8cm}@{}}
\toprule
Signal             & Width                         & Direction & Purpose \\
\midrule
\texttt{evict\_needed} & 1                         & out       & Asserted when SRAM is full \\
\texttt{evict\_addr}   & $\lceil\log_2 \texttt{SRAM\_DEPTH}\rceil$ & out & Address of entry to evict \\
\bottomrule
\end{tabular}
\end{table}

\noindent This interface connects to the Memory Hierarchy Controller (block~4).

%==============================================================================
% 4. Logical operations (compiler-facing)
%==============================================================================
\section{Logical Operations}
\label{sec:ops}

\begin{table}[h]
\centering
\small
\begin{tabular}{@{}l l l p{5.5cm}@{}}
\toprule
Operation                     & Inputs             & Outputs          & Description \\
\midrule
\texttt{KV\_QUERY}            & ---                & INFO struct      & Read all \texttt{INFO\_*} registers \\
\texttt{KV\_RESET}            & ---                & ---              & Soft reset: clear SRAM, reset FSM \\
\texttt{KV\_ENABLE}           & ---                & ---              & Set bit[1] of \texttt{CTRL} \\
\texttt{KV\_COMPRESS\_KEY}    & base, $G{\times}D$ coords & ---       & Buffer + compress a key \emph{group}; store $G$ per-token records \\
\texttt{KV\_COMPRESS\_VALUE}  & addr, $D$ coords   & ---              & Compress and store a value token at \texttt{addr} \\
\texttt{KV\_DECOMPRESS\_KEY}  & addr               & $D$ fp32 coords  & Read and decompress a key token from \texttt{addr} \\
\texttt{KV\_DECOMPRESS\_VALUE}& addr               & $D$ fp32 coords  & Read and decompress a value token from \texttt{addr} \\
\texttt{KV\_STATUS}           & ---                & STATUS bits      & Read the \texttt{STATUS} register \\
\texttt{KV\_OCCUPANCY}        & ---                & count            & Read the \texttt{OCCUPANCY} register \\
\bottomrule
\end{tabular}
\end{table}

\subsection{Worked Lowering Example}
\label{sec:lowering}

\textbf{Stage~1: setup.} Query hardware configuration once per compilation:

\begin{lstlisting}[basicstyle=\ttfamily\small]
info = KV_QUERY()   // -> dim=64, tier=1(CQ-4), group=128, outlier_k=0, sram_depth=16
KV_RESET(); KV_ENABLE()
\end{lstlisting}

\textbf{Stage~2: store phase.} Values are per-token; keys are grouped:

\begin{lstlisting}[basicstyle=\ttfamily\small]
for token_idx in new_tokens:                       // values: per token
    KV_COMPRESS_VALUE(token_idx % info.sram_depth, v_vectors[token_idx])

for group in key_groups_of(new_tokens, info.group): // keys: per G-token group
    KV_COMPRESS_KEY(group.base_addr, group.k_vectors)   // stream G*D coords
\end{lstlisting}

\textbf{Stage~3: attention phase.} Decompress cached K/V for attention:

\begin{lstlisting}[basicstyle=\ttfamily\small]
for cached_idx in range(num_cached):
    k_hat = KV_DECOMPRESS_KEY(cached_idx)   // D fp32 coords
    v_hat = KV_DECOMPRESS_VALUE(cached_idx)
    score = dot(query, k_hat) / sqrt(D)
    decision = PC_PUSH_TILE(score)          // block 1: precision controller
    out += (fp16_matmul if decision == FP16 else int8_matmul)(softmax(score), v_hat)
\end{lstlisting}

\subsection{Compiler Binding Patterns}
\label{sec:bindings}

\noindent\textbf{MLIR dialect.} Define \texttt{lhsi.kv.*} ops mirroring
\S\ref{sec:ops}; the lowering pass rewrites attention ops to emit
\texttt{lhsi.kv.compress\_key} (grouped), \texttt{lhsi.kv.decompress\_key}, etc.
\noindent\textbf{TVM Relax / TIR}, \textbf{ONNX CustomOp}, and \textbf{custom IR}
backends bind analogously to the C API in \S\ref{sec:capi}.

%==============================================================================
% 5. C API stub
%==============================================================================
\section{C API Stub}
\label{sec:capi}

\begin{lstlisting}[language=C]
/* lhsi_kv_cache_engine.h --- LonghornSilicon KV Cache Engine driver API (v0.2). */

#include <stdint.h>
#include <stdbool.h>
#include <stddef.h>

typedef struct {
    uint32_t vector_dim;       /* head dim D                                  */
    uint32_t tier;             /* 0=CQ-8, 1=CQ-4, 2=CQ-4+                      */
    uint32_t key_group;        /* G: tokens per per-channel key group         */
    uint32_t outlier_k;        /* top-k FP16 outlier key channels (CQ-4+)     */
    uint32_t scale_depth;      /* per-channel scale-bank depth (= D)          */
    uint32_t resid_depth;      /* residual-buffer depth (= G)                 */
    uint32_t sram_depth;
    uint32_t cr_k_fixed;       /* compression ratio K, fixed-point 8.8        */
    uint32_t cr_v_fixed;       /* compression ratio V, fixed-point 8.8        */
    uint32_t version;          /* 0x00020000 = v0.2                           */
} lhsi_kv_info_t;

typedef struct lhsi_kv_handle lhsi_kv_handle_t;

int  lhsi_kv_open (const char *device, lhsi_kv_handle_t **out);
void lhsi_kv_close(lhsi_kv_handle_t *h);
int  lhsi_kv_query(lhsi_kv_handle_t *h, lhsi_kv_info_t *out);
int  lhsi_kv_reset (lhsi_kv_handle_t *h);
int  lhsi_kv_enable(lhsi_kv_handle_t *h, bool enable);

/* Compress + store a KEY GROUP of `g` tokens (g <= key_group) at base_addr.
 * coords is g*dim FP16 values, token-major. */
int  lhsi_kv_compress_key_group(lhsi_kv_handle_t *h, uint32_t base_addr,
                                const uint16_t *coords_f16, size_t g, size_t dim);

/* Compress + store a single VALUE token at addr (dim FP16 values). */
int  lhsi_kv_compress_value(lhsi_kv_handle_t *h, uint32_t addr,
                            const uint16_t *coords_f16, size_t dim);

/* Read + decompress a key / value token: dim fp32 values out. */
int  lhsi_kv_decompress_key  (lhsi_kv_handle_t *h, uint32_t addr,
                              float *coords_f32, size_t dim);
int  lhsi_kv_decompress_value(lhsi_kv_handle_t *h, uint32_t addr,
                              float *coords_f32, size_t dim);

typedef struct { bool idle; bool sram_full; uint32_t occupancy; } lhsi_kv_status_t;
int  lhsi_kv_status(lhsi_kv_handle_t *h, lhsi_kv_status_t *out);
\end{lstlisting}

\noindent Return codes follow the standard convention: \texttt{0} on success,
\texttt{-errno} on failure.

%==============================================================================
% 6. Compression algorithm
%==============================================================================
\section{Compression Algorithm: ChannelQuant}
\label{sec:algorithm}

ChannelQuant is a per-channel/per-token integer quantization scheme designed for
KV-cache compression in GQA transformer inference. It follows the KIVI/KVQuant
recipe; the algorithm contract, reference model, and golden vectors are frozen in
the sibling \texttt{channelquant} repository (contract v0.2). Rounding is
round-half-to-even; INT4 clamps to $[-8,7]$, INT8 to $[-128,127]$; scales are
FP16 with $\text{EPS}=2^{-14}$.

\subsection{Key Path --- Per-Channel INT4}
Keys are quantized on the \emph{channel} axis, because the KV error on GQA models
is dominated by a few fixed high-magnitude key channels.
\begin{enumerate}
\item Buffer a group of \texttt{G} key tokens.
\item For each channel $c$, take $a_c = \max_t |k_{t,c}|$ over the group and form
      the FP16 scale $s_c = \max(a_c / q_{\max}, \text{EPS})$ with $q_{\max}=7$.
\item Quantize each keep channel: $\hat{q}_{t,c} =
      \mathrm{clamp}(\mathrm{round}(k_{t,c}/s_c))$ to INT4.
\item \textbf{CQ-4+ outlier lane}: the top-$k$ channels from the static calibrated
      ROM mask (\path{../channelquant/reference/masks/}, $k=2$) are stored FP16
      rather than INT4.
\end{enumerate}

\subsection{Value Path --- Per-Token INT4}
Each value token is scaled by $s_t=\max(\max_c|v_{t,c}|/q_{\max},\text{EPS})$ and
quantized to INT4 (INT8 with $q_{\max}=127$ for the CQ-8 tier). No grouping.

\subsection{Unified Per-Channel SRAM Record}
Each stored key token is a per-channel record
$\{\text{tag},\ D\times\text{FP16 field},\ D\times\text{INT4 code}\}$:
\begin{itemize}
\item \textbf{keep channel} $\to$ $\{$group scale $s_c$, INT4 code$\}$;
\item \textbf{outlier channel} $\to$ $\{$raw FP16 value, code $+1\}$, so that
      decompress $\text{code}\cdot\text{field}$ widens the FP16 exactly --- no
      separate sidecar region and no read-side mask.
\end{itemize}
Decompression is uniform per channel: $\hat{x}_c = \text{code}_c \cdot
\text{field}_c$ in fp32. Value records store $\{$FP16 scale, packed INT4/INT8
payload$\}$ and dequantize the same way.

\subsection{Effective Bits and Compression}
\begin{table}[h]
\centering
\small
\begin{tabular}{@{}l l l@{}}
\toprule
Path & Per-value bits & Notes \\
\midrule
Key (CQ-4)   & $4 + 16/G$        & INT4 code + amortized per-channel scale \\
Key (CQ-4+)  & $4 + 16/G + \frac{k}{D}(16{-}4)$ & plus the FP16 outlier lane \\
Value (INT4) & $4 + 16/D$        & INT4 code + per-token scale \\
\bottomrule
\end{tabular}
\end{table}
\noindent At $D{=}64$, $G{=}128$, $k{=}2$: mean $\approx$4.2 bits/value,
$\approx$3.8$\times$ vs.\ FP16's 16 bits.

\subsection{Serialized Datapath}
Each compute core (scale / quant / dequant) carries an FP16 divider. Rather than
$D$ parallel units (a single wide combinational divide cone that stalls
place-and-route), the datapath \textbf{serializes one shared unit} walked across
the $D$ channels, bit-exact with the behavioral oracle. This is what lets the
block synthesize, pass formal RTL$\equiv$netlist equivalence, and place-and-route
on Sky130 within the CI time budgets.

%==============================================================================
% 7. Synthesis-time configuration
%==============================================================================
\section{Synthesis-Time Configuration}
\label{sec:syntime}

\begin{table}[h]
\centering
\small
\begin{tabular}{@{}l c l p{5cm}@{}}
\toprule
Parameter             & Default & Range                   & Notes \\
\midrule
\texttt{VECTOR\_DIM}    & 64  & 16, 32, 64, 128          & Head dimension $D$ \\
\texttt{TIER}           & 1   & 0, 1, 2                  & 0=CQ-8, 1=CQ-4, 2=CQ-4+ \\
\texttt{KEY\_GROUP}     & 128 & $\geq 2$                 & Tokens per per-channel key group ($G$) \\
\texttt{OUTLIER\_K}     & 0   & 0, 2                     & FP16 outlier key channels (CQ-4+) \\
\texttt{SCALE\_WIDTH}   & 16  & 16                       & FP16 scale width \\
\texttt{COORD\_WIDTH}   & 16  & 16                       & FP16 input coordinate width \\
\texttt{OUT\_WIDTH}     & 32  & 32                       & fp32 decompressed output width \\
\texttt{SRAM\_DEPTH}    & 16  & 2--4096                  & Behavioral reg array for Sky130; real macros at 16FFC \\
\bottomrule
\end{tabular}
\end{table}

\noindent\textbf{Gate proxy.} The synthesis / formal / OpenLane CI gates run a
deliberately small default proxy (the SRAM and residual buffer are behavioral
flip-flops, no Sky130 macro); the shipped $D$/$G$/depth are set per instantiation.

%==============================================================================
% 8. Bit-Accurate Reference
%==============================================================================
\section{Bit-Accurate Reference}
\label{sec:ref}

Three implementations are bit-exact with each other (3-way parity, contract \S8):
\begin{enumerate}
\item \textbf{Python}: \path{../channelquant/reference/channelquant_ref.py} ---
      the frozen algorithm reference + golden vectors.
\item \textbf{C++}: \path{sw/reference_model/channelquant_ref.{hpp,cpp}} ---
      a 1:1 port of the RTL behavioral cores.
\item \textbf{SystemVerilog}: \path{rtl/kv_cache_engine.sv} + \path{cq_key_path.sv}
      / \path{cq_value_path.sv} / \path{cq_units_syn.sv} + support --- synthesizable RTL.
\end{enumerate}

\noindent The RTL is verified bit-exact against the golden vectors by
\path{make sim_cq} / \path{sim_kpath} / \path{sim_vpath} / \path{sim_top}, and the
C++/Python legs by \path{make -C sw/reference_model test-all}.

%==============================================================================
% 9. Integration phases
%==============================================================================
\section{Integration Phases}
\label{sec:phases}

\begin{table}[ht]
\centering
\small
\renewcommand{\arraystretch}{1.2}
\begin{tabular}{@{}c p{3.0cm} p{4.0cm} p{5.0cm}@{}}
\toprule
Phase & Timeline & Compiler targets & We provide \\
\midrule
0 & now & Bit-accurate Python and C++ reference models & Models + ISA + tests + Sky130 sign-off \\
1 & when FPGA arrives & AXI-Lite + AXI-Stream on Zynq UltraScale+ & Vivado bitstream + driver \\
2 & after all 4 blocks complete & Multi-block FPGA project & Integrated attention pipeline \\
3 & post-tape-out (2027+) & PCIe-attached custom chip & TSMC 16FFC silicon + Linux driver \\
\bottomrule
\end{tabular}
\end{table}

%==============================================================================
% 10. Interaction with other blocks
%==============================================================================
\section{Interaction with Other Blocks}
\label{sec:interaction}

\begin{itemize}
\item \textbf{Block 1 (Precision Controller / ACU)}: receives the fp32
      decompressed K/V from this block for attention; it decides INT8 vs FP16
      routing per tile from the scores. The controller is codec-agnostic --- it
      never sees the KV storage format --- so the ChannelQuant codec change
      requires no ACU change.
\item \textbf{Block 3 (Token Importance Unit)}: provides token-importance scores
      informing eviction when SRAM is full.
\item \textbf{Block 4 (Memory Hierarchy Controller)}: receives
      \texttt{evict\_needed} / \texttt{evict\_addr} when SRAM is full.
\end{itemize}

%==============================================================================
% 11. Open questions
%==============================================================================
\section{Open Questions}
\label{sec:open}

\begin{enumerate}
\item Partial-group flush ($g<G$): the datapath supports it; the top's stream
      framing currently auto-flushes at full $G$ only. Which sideband signals a
      partial group (control-register bit vs.\ stream marker)?
\item Should the outlier count $k$ and the ROM mask be run-time reloadable
      (per deployed model) rather than fixed at synthesis?
\item Is a per-token DWB-style tier selector (2/4/8-bit) worth the controller
      cost, given per-channel INT4 already recovers $\sim$FP16?
\item Are the fp16 scale / fp32 read-bus widths sufficient for larger heads
      ($D=256$)?
\item Should the eviction policy be configurable via registers (LRU, FIFO,
      importance-based)?
\end{enumerate}

%==============================================================================
% 12. Change log
%==============================================================================
\section{Change Log}
\label{sec:changelog}

\begin{itemize}
\item \texttt{kv-isa-0.2} (2026-07-03): Codec replaced TurboQuant+ $\to$
      ChannelQuant (per-channel INT4 keys, per-token INT4 values, CQ-4+ outlier
      lane). New \texttt{INFO\_TIER}/\texttt{GROUP}/\texttt{OUTLIER\_K}/
      \texttt{SCALE\_DEPTH}/\texttt{RESID\_DEPTH} registers; \texttt{INFO\_PQ\_BITS}
      / \texttt{INFO\_QJL\_BITS} removed. fp32 read bus. RTL fully integrated and
      signed off (all CI gates green, Sky130).
\item \texttt{kv-isa-0.1} (2026-05-14): First public draft (TurboQuant+ codec,
      retired).
\end{itemize}

\vspace{0.5in}
\noindent\rule{\textwidth}{0.4pt}\\
\noindent\small\textit{LonghornSilicon KV Cache Engine --- Interface Specification.
This document is open and may be redistributed. The TSMC 16FFC PDK referenced
in the integration phases is NDA-protected and is not part of this document
or the public repository.}

\end{document}
